\uuid{HB0z}
\exo7id{2760}
\titre{exo7 2760}
\auteur{tumpach}
\organisation{exo7}
\datecreate{2009-10-25}
\isIndication{false}
\isCorrection{false}
\chapitre{Déterminant, système linéaire}
\sousChapitre{Calcul de déterminants}
\cpge{mpsi-mp2i.determinants.d, pcsi.calcul-matriciel-systemes-lineaires.d, ptsi.calcul-matriciel-systemes-lineaires.d, mpsi-mp2i.calcul-matriciel-systemes.e, pcsi.determinants.c, ptsi.determinants.a, mpsi-mp2i.calcul-matriciel-systemes.c, pcsi.calcul-matriciel-systemes-lineaires.c, ptsi.calcul-matriciel-systemes-lineaires.c, bcpst1.systemes-lineaires.a, bcpst1.matrices}
\module{Algèbre}
\niveau{L2}
\difficulte{}

\contenu{
\question{
En utilisant le d\'eterminant montrer que chacun des syst\`emes suivants admet une solution unique.  R\'esoudre chacun de ces syst\`eme en inversant la matrice de ses  coefficients~:  
$$
\left\{\begin{array}{ccccccc}
x & + & y & + & z & = &1\\
2x & +  & 3y & + & 4z &  =  & 2\\
&& y &+ &  4z & = & 3 
\end{array}\right.
$$
$$
\left\{\begin{array}{cccccccc}
x & + & y & + &  z & + &  t  & = 1\\
x & + &y & - & z & - & t & =  0\\
x & - & y & - & z&  +  &t & = 2\\
x & - & y & + & z & - & t & = 3  
\end{array}\right.
$$
\vspace{.2cm}
$$
\left\{
\begin{array}{cccccc}
3x & + & 2y & - &  z &  = -1\\
x&+&y&-&3z& = 1\\
3x& + &2y &&&= 0\\
\end{array}
\right.
$$
}
}
